\section{Dependencies and open questions}\label{sec:discussion}
The quantitative scalar-to-support-to-rounding chain is written out in Sections~\ref{sec:core}--\ref{sec:global}. Its parity-free part proves the all-half-order construction, while the half-integer trace is the only oddness input in the global odd theorem. The graph/kernel equivalence and its local conditioning are independent of global stability; the residual-only entrance estimate uses the global theorem. Component obstructions use graph pairings and elementary spectral and binary identities. The quadratic-family argument uses direct counts, the graph bound, and the global entrance estimate, not the two small rank computations. The second-order conjugate-power proof is independent of quantitative stability, and the local exponent classification combines these separate ingredients with exact finite-circle geometry.

Several questions remain unresolved.
\begin{enumerate}
\item Is the rectangle graph connected for every cyclic $\GH(m,2)$ with odd $m$? Neither a proof nor a disconnected odd seed is supplied. A proof would yield the polynomial linear estimate at every odd seed; a counterexample would give sharp square-root behavior there.
\item Which abstract component-weight patterns are realized? The exclusions leave patterns such as $4+6$ and $5+5$ at $m=5$; they are not a complete classification.
\item No new-order seed existence, including unresolved cyclic prime-power or non-prime-power existence questions, is asserted. At even half-orders, determining which seeds admit compatible rectangles remains a separate realization problem.
\item The ordinary first-power isolation conjecture and zero ordinary defect for the whole quadratic family remain outside these results.
\item Constants and dependence on $m$ have not been optimized. The inverse-minor estimate is not polynomial, while the connected-graph estimates are.
\item Historical comparison is incomplete. In particular a full-text comparison with Craigen--Woodford~\cite{CW} remains necessary before any novelty claim.
\end{enumerate}
The terminology ``almost Hadamard'' in~\cite{BNZ} refers to real matrices whose normalization is orthogonal and locally maximizes the entrywise $1$-norm; that is a different problem from the present all-power residual. S-Hadamard matrices and square-defect questions~\cite{Phangara} are related but require fewer powers than the full range when $m>3$. Neither comparison justifies a priority claim. The ordinary defect framework remains that of~\cite{TZ}.
